% =====================================================================
% Related Literature Review, v7
% DRAFT (Claude, 2026-11-10, with Anders). Built on the structure and text
% of "Literature_review_pumpekraft_revidert (2).pdf" (5.1-5.7), which Anders
% likes, and deepened along the assessment he pasted (16 discussion points
% from the original draft, more evidence on findings, effect sizes where
% they carry the argument, physical detail in the methods section, price
% volatility and persistence developed before the gap). Claims are checked
% against the paper notes (literature/papers/<bibkey>/note.md, which give
% page numbers), not against paper.txt directly. Weaknesses of the sources
% are taken from the notes. See drafts/literature-review-v7/CHANGES.md.
% Citations: natbib commands through biblatex natbib=true; keys from
% ../pumped-hydro-thesis/references.bib.
% =====================================================================

\section{Related Literature Review}
\label{sec:literature}

The literature on pumped-storage hydropower addresses several related
questions: whose value is assessed, which uncertainties affect operation,
which markets provide revenue, and how the plant's capacities affect its
value. This chapter reviews these questions in separate thematic
subsections, in which the studies are compared, contrasted and critiqued,
and connected to the present study. Section~\ref{sec:lit_scope} defines
the scope, Section~\ref{sec:lit_perspective} examines valuation
perspectives, Section~\ref{sec:lit_uncertainty} discusses uncertainty and
price behaviour, Section~\ref{sec:lit_markets} examines market
participation, and Section~\ref{sec:lit_capacities} considers capacity
configurations. Section~\ref{sec:lit_methods} then compares the solution
approaches used in the literature and the physical detail they retain.
Section~\ref{sec:lit_gap} draws the findings together and identifies the
research question motivating the present study.

% ---------------------------------------------------------------------
\subsection{Scope and Selection of Literature}
\label{sec:lit_scope}

The review focuses on operational scheduling and economic valuation of
pumped-storage hydropower. Particular attention is given to studies that
connect short-term operating opportunities with longer-term reservoir
management. Studies of conventional reservoir hydropower are included
where they offer relevant insights into water values, uncertainty or
reserve markets. System-planning and hybrid-energy studies provide
context for evaluating storage under different objectives, and studies
of electricity-price behaviour are included because the research question
concerns price volatility and persistence.

\todonote{TODO (authors): Describe the search process you actually used:
databases (e.g.\ Google Scholar, Scopus), search terms, inclusion
criteria (peer-reviewed, relevance of title and abstract), and whether
further studies were found by following reference lists (backward
snowballing). Bore and Brensholm describe how they extended the review of
their project report by backward snowballing. Three to five sentences
are enough. The search log (\texttt{literature/search-log.md}) is still
empty.}

The chapter provides a focused narrative synthesis of selected studies.
Their relevance is assessed through the economic question addressed, the
representation of operation, the market opportunities considered and the
treatment of plant capacities. Individual studies may appear in several
subsections, but each discussion concerns the dimension relevant to that
subsection. This allows the literature to be compared thematically while
distinguishing findings that arise under different assumptions.

% ---------------------------------------------------------------------
\subsection{Valuation Perspective}
\label{sec:lit_perspective}

The perspective adopted in a study determines what is meant by the value
of pumped storage. System-oriented studies assess its contribution to
the wider electricity system, whereas owner-oriented studies assess the
economic benefit available to an individual producer. The two answer
different questions, and conclusions cannot be transferred directly
between them.

\citet{helseth2013} examine pumped storage within a hydrothermal system
that also includes wind generation and transmission constraints, and
\citet{helseth2017} assess a pumped-storage expansion through changes in
expected socio-economic surplus. \citet{zhang2026unlocking} compare
seasonal and conventional pumped storage in a Chinese system with much
wind and solar power. The seasonal plant reduces curtailment, emissions
and operating cost, but at higher capital cost and a higher levelised
cost of electricity. Such results show what storage contributes to the
system, but a favourable system-level assessment does not establish the
corresponding private return, since benefits can also accrue to
consumers and other producers. Likewise, additional owner revenue does
not measure all the system benefits.

Owner-oriented studies instead evaluate the income associated with
operating a particular asset. \citet{ak2019quantifying} report revenue
gains from operating an existing hydropower cascade in pumped-storage
mode, and \citet{alic2023optimal} assess the economic benefit of adding a
pump to a Norwegian cascade. Within an owner perspective, the treatment
of market power introduces a further distinction.
\citet{abgottspon2013stochastic} consider a price-making producer whose
trading decisions influence prices, whereas \citet{helseth2016} consider
a price-taking producer facing given prices. A price taker responds to
available price differences, while a price maker also accounts for how
its own actions reduce them.

Increased operating revenue alone does not establish investment
profitability. \citet{bore2026assessing} take this further step and use
additional revenues to estimate the maximum investment cost that a
capacity upgrade can support, which requires an economic lifetime, a
discount rate and cost assumptions. \citet{ak2019quantifying} make a
simpler check by comparing the annual revenue gain with the annuity of
the extra investment in reversible turbines. Conclusions about the
economic value of storage must therefore be read against the objective
of the study from which they are drawn.

% ---------------------------------------------------------------------
\subsection{Uncertainty in Pumped-Storage Scheduling}
\label{sec:lit_uncertainty}

Uncertainty affects the trade-off between using stored water now and
retaining it for future operation. For pumped storage, future prices
influence both the opportunity to sell electricity and the cost of
replenishing storage through pumping. Natural inflows provide an
additional uncertain source of water in inflow-dependent plants. The
literature differs in which of these uncertainties is represented, how
they are described over time, and what information is available when
operating decisions are made.

\subsubsection{Sources of Uncertainty and Their Dependence}
\label{sec:lit_unc_sources}

Electricity prices are central to the owner-oriented studies reviewed
here. \citet{alic2023optimal} account for uncertainty in both prices and
inflows in a pumped-storage cascade, and \citet{bore2026assessing}
likewise consider uncertain inflows and prices for conventional
reservoir hydropower. \citet{helseth2016} additionally treat
reserve-capacity prices as uncertain. Other uncertain inputs become
relevant when pumped storage is studied as part of a wider energy
system. \citet{naval2025optimization} consider wind uncertainty in the
operation of a hybrid system, whereas \citet{zisos2023mixing} examine
uncertainty in wind, demand and energy conversion in an island-energy
design problem. In both, prices play a subordinate role: the first
treats a full year of day-ahead prices as known, and the second contains
no prices at all. Inflow assumptions appropriate to a reservoir-fed
cascade also need to be distinguished from those of a closed-loop plant
without inflow, such as the one studied by \citet{chazarra2017}.

When several inputs are uncertain, their dependence becomes a modelling
choice. \citet{kleiven2023} find that ignoring the negative correlation
between Norwegian prices and inflow lowers water values and leads to a
higher reservoir in autumn, but changes expected revenue very little. \citet{helseth2016} treat prices as independent of inflow
and justify this only qualitatively, and \citet{perezdiaz2020medium} find
the correlation between energy and reserve prices too weak to affect
water values for a Spanish plant. These findings suggest that dependence
matters more for how the reservoir is operated than for the revenue it
earns. They concern conventional plants, however, and do not show
whether the same holds when prices also drive pumping.

\subsubsection{Timing and Resolution of Information}
\label{sec:lit_unc_timing}

Long-term storage decisions depend on the short-term operating
opportunities available when the water is eventually used. Seasonal
inflows and prices govern how the reservoir is managed over the year,
whereas price variations within the day and week create the
opportunities for pumping and generation. \citet{ak2019quantifying}
illustrate the division: when their cascade is operated as pumped
storage, the monthly reservoir decisions barely change, while the
allocation of pumping and generation over the hours of the day changes
considerably. A long evaluation horizon therefore does not by itself
establish that a study captures storage flexibility.

The timing of information distinguishes studies even when they include
the same uncertain variables. \citet{bore2026assessing} represent
uncertain conditions between weeks but assume that the current week's
inputs are known when its decisions are made. Their model therefore
combines future uncertainty with perfect information within each week.
An hourly operating schedule does not, by itself, imply that uncertainty
is resolved hour by hour. \citet{helseth2013} make the same weekly
assumption and note themselves that it is too coarse to capture the
balancing needs created by wind power.

\citet{abgottspon2014} examine when this simplification matters. For a
Swiss pumped-storage plant, they compare four ways of representing
operation within the week, from an aggregated model to one in which
information arrives within the week. Treating the week as known
overvalues water. A stochastic representation of the week gives about
10\% higher expected profit, and a similar improvement in the poor
outcomes, but only when weekly secondary reserves are included; without
reserves, the model with a known week performs better. The authors
conclude that stochastic within-week detail is worthwhile when the
market structure is sufficiently complex, and suggest that complex plant
structures may be another reason, but leave this untested. The result
rests on a single plant whose capacities are not reported, and profits
are evaluated with an operator's dispatch heuristic, so differences
between the models mix policy quality with the heuristic's behaviour.

Time resolution and information are related but separate choices.
\citet{abgottspon2016multihorizon} show that, for a Swiss system with
seasonal storage reservoirs and small balancing reservoirs, a purely daily model undervalues
and a purely hourly model overvalues the plant by more than 10\%,
whereas linking daily and hourly decisions comes within a few per cent.
Their reported case is deterministic, however.
\citet{alic2023optimal} find that simulating operation day by day rather
than week by week lowers revenue considerably in weeks when a
state-dependent environmental restriction becomes active, because the
weekly model knows the reservoir state too coarsely. Testing several
price scenarios can establish sensitivity to different conditions, but
it does not by itself establish how decisions adapt as information
arrives.

\subsubsection{Representing Uncertainty over Time}
\label{sec:lit_unc_repr}

Uncertainty is typically represented by a small number of weekly states
derived from data, and the reviewed studies share several steps.
Historical or simulated series are collected, grouped into a few typical
states per week, and transition probabilities between states are
estimated. \citet{alic2023optimal} cluster weekly prices and inflows
into five joint states, and \citet{bore2026assessing} into three.
\citet{helseth2016} build nine states from three energy-price and three
reserve-price clusters of only 80 simulated scenarios, and have to adjust
the transition probabilities slightly because so few scenarios are
available for nine states. \citet{mo2001} describe how such probabilities can
be corrected so that the conditional expected prices are preserved. The
number of states is thus constrained by the data: more states describe
the variation better but rest on fewer observations each.

Aggregation within the state affects what the pump can earn.
\citet{alic2023optimal} represent each state by weekly average prices,
spread over the week with a fixed daily profile, so the price variation
within the week is the same in every state. \citet{bore2026assessing}
use the median series of each cluster, and \citet{ak2019quantifying}
represent hourly prices through representative daily price curves for
each month, which removes the chronology of pumping and generation. For
pumped storage, whose income comes largely from short price differences,
such averaging may understate the value of pumping. This may help
explain why \citet{alic2023optimal} find a small gain from their pump,
although the authors do not make this argument. An alternative is to use
historical years directly. \citet{helseth2017} let each historical
weather and price year occur with equal probability and correct it to
the current state with a simple autoregressive model, and
\citet{helseth2013} simulate operation on historical records to retain
correlations that their stochastic models omit. Data granularity should
therefore be chosen with the source of the plant's income in mind.

\subsubsection{Price Volatility and Persistence}
\label{sec:lit_unc_price}

A statement that a study ``includes uncertainty'' is not sufficient to
characterise it. For a storage owner trading on price differences, the
representation of price behaviour over time is likely to be at least as
important as the number of uncertain inputs. Two properties of this
behaviour are central: volatility, how much prices vary around their
expected pattern, and persistence, how long high- or low-price
conditions last (Section~\ref{sec:ph25-volatility-persistence}).

\citet{haldrup2006regime} find that Nordic day-ahead prices are highly
persistent, most of all in the Norwegian price areas, which they link to
the slow variation of reservoirs and inflow. Their data are from
2000--2003, however, before market coupling and the growth of wind and
solar power, and their regimes last hours rather than weeks. A weekly
state description captures persistence only through its transition
probabilities, and \citet{mo2001} note that such a description forgets
its initial state at an exponential rate, so uncertainty about prices far
ahead needs an additional model. Future price behaviour may also change:
\citet{henden2016economic} find that more transmission to the continent
increases the variability of southern Norwegian prices, whereas more
pumped storage reduces it.

Studies that relate price behaviour to storage value do so mostly with
known prices. \citet{ak2019quantifying} find the largest gain from
pumping in the price year with the largest variation within the day,
although the relation is not one-to-one. \citet{bradbury2014economic}
find that higher volatility raises the return on storage in seven US
markets. \citet{sioshansi2009estimating} show that a simple rule based
on the prices of the previous two weeks captures most of the
perfect-foresight arbitrage value, because daily and weekly price
patterns repeat. Short-term variation also matters more for pumped
storage than for conventional hydropower: \citet{lohndorf2023} find that
coordinating day-ahead and intraday trading is worth about an order of
magnitude more for a pumped-storage plant than for a reservoir plant
without a pump, consistent with the small value of intraday trading that
\citet{faria2011day} find for a conventional cascade. None of these
studies varies volatility and persistence separately or relates them to
the value of pump, turbine and reservoir capacity under uncertainty.
Since pumping and generating power earn from short price differences,
whereas reservoir capacity earns from differences that last over weeks
or seasons, the two properties may favour different capacities.

% ---------------------------------------------------------------------
\subsection{Electricity Market Participation}
\label{sec:lit_markets}

Market participation determines the revenue opportunities available to
a storage owner. Energy-market participation rewards electricity sales
relative to the cost of purchases for pumping. Other products can
remunerate the availability of operating flexibility or its actual use.
Studies must therefore be compared by the products and settlements
included, since these define what their reported revenues cover. The
reported gains are not directly comparable either, as market design,
price levels, plants and comparison baselines differ between the
studies.

\paragraph{Energy and forward trading}
Energy-centred studies provide a baseline for assessing the value of
pumping. \citet{ak2019quantifying} and \citet{alic2023optimal} value
pumping against day-ahead prices only, and \citet{alic2023optimal} add
environmental restrictions that limit how fast the reservoir level may
change. \citet{abgottspon2013stochastic} add forward contracts for a
price-making Swiss producer. Physical production and trading positions
are then separated: in one simulated week, a forward obligation is met
by buying cheaper energy back in the day-ahead market rather than by
generating. Forward trading raises simulated annual profit by about
20\%. The result depends on the price-making assumption, with a price
response estimated from experience, and although the study is described
as stochastic in prices, only inflow is represented as random. Forward
positions extend the producer's commercial opportunities, but they serve
a different purpose from payments for maintaining reserve availability.

\paragraph{Reserve capacity}
Reserve participation changes scheduling and water values as well as
revenue. \citet{helseth2016} compare energy-only scheduling with joint
sales of energy and symmetric reserve capacity, traded the week before
energy, for a Norwegian watercourse of seven plants. A large part of the
reserve income is earned without changing the generation schedule,
because the plants are most efficient below full output and therefore
have spare capacity in both directions. Where the schedule does change,
the large reservoirs are kept higher until the spring flood and then
used more to keep generators running in the low-load season, some of
them at minimum output. In their case, water values are lower with reserves, since
capacity is withheld from generation in periods of high output.
\citet{helseth2017iet} find a similar pattern for a single reservoir in
southern Norway: selling primary reserves raises expected operating
profit by 7 to 11\%, shifts less energy between seasons, and makes water
values fall faster as the reservoir fills. Both studies evaluate their
policies on scenarios drawn from the same model used to compute them,
and the second fits its price model to only three years of data.

Reserve commitments compete with other uses of the same flexibility. In
\citet{helseth2017iet}, energy revenue falls when reserves are sold, so
the net gain is smaller than the reserve income. \citet{chazarra2017}
find that a closed-loop plant earns most of its income from secondary
reserve capacity, and that the most flexible configuration accepts a
loss in the day-ahead market in order to offer reserves in every hour.
\citet{perezdiaz2020medium} report a gain of about 4.5\% from including
the automatic frequency restoration reserve market when computing water
values for a Spanish plant, and find that accounting for its own effect
on reserve prices reduces spillage considerably while barely changing
profit. \citet{bore2026assessing} find that capacity markets make up a
larger share of revenue in periods of high inflow and low energy prices,
and that attractive capacity prices can make the plant spill water and
restrain production to keep capacity available for upward reserves.
Reserve income must therefore be evaluated together with the operating
opportunities given up to provide it.

A pump widens what a plant can offer. A conventional plant cannot
provide reserves while pumping, which is why the configurations of
\citet{chazarra2017} that can pump and generate at the same time earn
more. \citet{dallinger2019socioeconomic} show that adding pump-turbines
to an Austrian reservoir plant extends its operating range to between
about $-70$\% and 100\% of turbine capacity, so it can offer downward
reserve without generating and no longer depends on natural inflow to
deliver reserves. \citet{abgottspon2014} include weekly secondary
reserves for a Swiss pumped-storage plant and find that they raise
expected profit in all four of their models.

\paragraph{Reserve activation}
The treatment of activation further distinguishes market
representations. \citet{chazarra2017} include regulation energy
associated with upward and downward reserve use, based on historical
activation, and show that storage limits restrict reserve offers: the
conventional plant does not offer reserves in an hour in which the
expected activation could empty its upper reservoir. Their results rest
on two representative weeks of Spanish data with prices and activation
known in advance. \citet{dallinger2019socioeconomic} simulate
activations every quarter hour but offer capacity at a price of zero, so
their reserve revenue comes from activated energy only.
\citet{bore2026assessing} include day-ahead and intraday energy alongside
FCR-N, FCR-D and upward and downward aFRR capacity, but exclude
activation. Compared with \citet{chazarra2017}, this shows that the
number of included products and the completeness of their delivery
representation are separate dimensions of market coverage.

Consequently, studies of market participation need to distinguish energy
trading, reserve-capacity commitments and reserve activation. Revenues
from these activities interact because they depend on shared operating
capability. Findings from different market combinations cannot be
compared solely through their total income, particularly when their
treatment of activation and settlement differs. Most of the multi-market
evidence also concerns conventional reservoir hydropower. For pumped
storage, only \citet{abgottspon2014} combine reserves with the long-term
value of water, for a plant whose capacities are not reported.

% ---------------------------------------------------------------------
\subsection{Pumping, Generating and Reservoir Capacities}
\label{sec:lit_capacities}

Capacity studies differ in which plant characteristics are fixed and
which are allowed to change. Optimising the operation of an existing
plant, comparing specified upgrades and optimising the plant's
dimensions address different questions. An effective operating policy
for one configuration does not establish whether that configuration is
well sized.

\paragraph{Adding a pump}
The most concrete configuration comparison in the literature is the
addition of a pump to an existing cascade, and the reported gains are
modest. \citet{ak2019quantifying} find that operating a Turkish cascade
as pumped storage raises turbine revenue by a quarter to a third, but
net revenue by only about 6\% on average after the cost of pumping. Total
net energy supplied falls, most of the gain comes from the reservoir with
the largest storage, and the most downstream plant loses revenue because
no water is pumped into it. \citet{alic2023optimal} add a
\SI{38}{\mega\watt} pump to a Norwegian cascade and find that median
energy production rises by about 9\% but median revenue by only about
2\%. The production figure counts gross turbine output, so recirculated
water is counted twice, and only the revenue reflects the cost of
pumping. Their pump also has the same power as the existing turbine, so
the effect of pump size cannot be separated from that of adding a pump
at all. \citet{schuller2024practical} value a new plant with pumping in
the Sira-Kvina system and find that the gain is smaller when operation is
simulated with hourly detail than in the medium-term model, because
short-term restrictions such as a minimum head for pumping limit the use
of the new plant. All three studies use the energy market only.

\paragraph{Upgrading generating units}
\citet{helseth2017iet} compare two reinvestment alternatives of the same
total generating power. The configuration with two Pelton units gains
more from reserve sales than the one with a large Francis unit, because
Pelton units can run at low output and leave more room for reserves,
although the Francis unit has the higher peak efficiency. The value of a
configuration thus depends on the markets included.
\citet{bore2026assessing} examine generating-capacity upgrades of 30, 60
and \SI{90}{\mega\watt} for a conventional plant. Energy revenue hardly
changes with the upgrade, because the plant has no more water to sell,
while capacity-market revenue grows substantially, and the marginal
benefit of each step falls. Additional generating power is thus limited
by the water available, and its value depends on markets that pay for
capacity rather than energy. Their alternatives concern generating power
only, in a plant without a pump.

\paragraph{Large expansions and market conditions}
Large capacity additions can change the prices they rely on.
\citet{helseth2017} expand the \SI{200}{\mega\watt} Duge pumped-storage
plant by up to \SI{2000}{\mega\watt} in a model of the Nordic power
system, in which prices follow from the system operation. A detailed
representation of the watercourse gives a lower marginal value of the
additional capacity than an aggregated model, because the detailed
constraints of the surrounding reservoirs limit its use. The paper does
not state whether pump and turbine are expanded separately, and the
pumped-storage case relies on an assumed exchange capacity to Germany
chosen to create price differences that trigger pumping.
\citet{henden2016economic} find that a large expansion of pumped storage
in Norway removes much of the price differences it trades on and leaves
its owners with little income, and \citet{dallinger2019socioeconomic}
find that the price spread captured by Austrian pumped storage falls as
more storage is built. An expansion of this scale is therefore valued
differently from an upgrade evaluated under given prices.

\paragraph{Sizing the reservoir}
Design optimisation brings selected dimensions into the decision problem
itself. \citet{zisos2023mixing} optimise the upper reservoir and the
photovoltaic capacity of an island system while holding pumping and
generating power fixed at \SI{6}{\mega\watt} each. Their benchmark
optimum is a reservoir of about \SI{315000}{\cubic\metre}, less than a
third of an earlier proposal of \SI{1.1}{\mega\cubic\metre}, although
the authors caution that the two rest on different assumptions. Across
weather and demand scenarios, the scenarios that call for a deeper
reservoir give lower profit, and the authors use this relation to show
the range of profit a designer accepts by choosing a larger, more
conservative reservoir. The plant follows a fixed rule that pumps all
surplus and covers all deficits, without prices, and each scenario has
its own optimal design, so the study does not identify one design that
performs well across scenarios. Evidence from electricity markets points
to diminishing returns from additional storage duration.
\citet{sioshansi2009estimating} find that most of the arbitrage value of
a storage plant is earned by its first hours of storage, and
\citet{bradbury2014economic} find that the most profitable ratio of
storage to power is set mainly by efficiency and by how costs divide
between power and storage. Both studies use known prices and only
energy arbitrage. \citet{chazarra2017} further observe that their most
flexible configuration uses only a small part of its reservoir, and
suggest, without testing it, that this could affect reservoir sizing.

The reviewed studies therefore provide evidence on operating given
assets, introducing pumping, comparing expansions and selecting
reservoir dimensions. Together they indicate what limits the value of a
capacity: the water available, the cost of pumping, short-term
operating restrictions, the markets that pay for flexibility, and the
effect of the plant itself on prices. They differ in which capacities
are varied and which remain fixed. This matters because an additional
unit of one capacity may have limited value when another capacity is
binding; a larger pump, for example, offers little additional storage
opportunity when the reservoir is already full. None of the reviewed
studies varies pumping power, generating power and reservoir capacity
jointly for a plant that trades in electricity markets.

% ---------------------------------------------------------------------
\subsection{Solution Approaches for Stochastic Hydropower Scheduling}
\label{sec:lit_methods}

The reviewed studies use different approaches to connect current
operation with uncertain future opportunities. Comparing them requires
distinguishing the formulation of the decision problem, the algorithm
used to solve it, and the procedure for updating decisions over time.
The principles are introduced in
Sections~\ref{subsec:two_stage_stochastic_optimisation},
\ref{subsec:stochastic_dynamic_programming} and
\ref{subsec:stochastic_dual_dynamic_programming}; this section compares
how they have been applied and what each application had to give up.

\paragraph{Physical representation}
Physical operating detail affects the mathematical problem to be solved,
and the studies differ widely in how they represent pumping, generation,
head and efficiency. \citet{chazarra2017} use a mixed-integer
formulation to represent the operating modes of each unit, including
hydraulic short-circuit operation, in which water is pumped and
generated at the same time. With ternary units this mode more than
doubles the income of their closed-loop plant, whereas with reversible
units the gain is small. \citet{ak2019quantifying} use a nonlinear model with
head-dependent generation but the same efficiency curve for pumping and
generation. \citet{helseth2017iet} represent head-dependent
generation curves and unit commitment with integer variables. Several
studies instead assume constant head and efficiency
\citep{alic2023optimal,naval2025optimization,zisos2023mixing}, and
\citet{helseth2017} treat head as a parameter updated each week. These
choices determine which operating behaviour a model permits: a
continuous model allows low output and small reserve adjustments that
real units may not deliver, and a constant efficiency makes every
pumping cycle equally profitable. Their effect on conclusions has been
measured for profit, but not for the ranking of plant configurations.

\subsubsection{Two-Stage Stochastic Programming}

Two-stage formulations distinguish decisions taken before uncertainty is
resolved from subsequent decisions that depend on its realisation. A
second-stage scenario can cover many operating periods, even though it
represents only one further information stage. \citet{faria2011day} use
a two-stage stochastic mixed-integer programme for the day-ahead bidding
of a Norwegian cascade, with intraday trading as the second stage, and
find that anticipating intraday trading changes profit very little. \citet{helseth2017} use two-stage stochastic linear problems for
detailed hydrothermal scheduling of the Nordic system, with scenarios
for weather and exogenous prices resolved in the second stage and solved
by Benders decomposition. This retains the full hydropower system
without aggregation, but simplifies the future information structure:
each scenario is known once the first week has passed, so the plan for
later weeks is made with foresight. Head is represented as a parameter
with precomputed sensitivities rather than as a variable, so the linear
formulation does not capture non-convex operating relationships.

\subsubsection{Rolling-Horizon Planning}

Rolling-horizon planning updates an operating plan repeatedly as new
information becomes available. The initial decisions are implemented,
the system state is updated, and the problem is solved again. This
procedure can be applied to deterministic or stochastic formulations.
\citet{helseth2017} embed their two-stage problems in a rolling-horizon
simulator that moves one week at a time. To keep this tractable, they
reduce the number of historical scenarios, use a coarser time resolution
further into the future, and take the value of water at the end of the
scenarios from another model; even so, a simulation takes 50 to 156
hours. Deterministic rolling schemes are cheaper.
\citet{sioshansi2009estimating} optimise two weeks at a time with a
15-day horizon so that stored energy keeps a value at the end of each
window, and \citet{dallinger2019socioeconomic} optimise each week with a
two-week look-ahead. Re-optimisation introduces adaptation between
successive solves, while each individual solve retains its assumed
information structure. Such approaches should therefore be compared
with a multistage formulation on the basis of how decisions can adapt,
rather than on the total simulation horizon alone.

\subsubsection{Stochastic Dynamic Programming}

Stochastic dynamic programming (SDP) connects current storage with
future operating opportunities through a value function computed on a
grid of reservoir levels, and places no convexity requirement on the
stage problems. \citet{alic2023optimal} choose SDP because their plant
is subject to a state-dependent water-abstraction restriction, which is
non-convex, alongside minimum environmental flows and limits on how fast
the reservoir level may change. The restrictions matter: the ramping
limits move production from high-price to lower-price periods and reduce
revenue, and a finer grid improves accuracy at a computational cost of
several hours for two reservoirs. \citet{helseth2017iet} combine SDP with
weekly mixed-integer problems. A linear version of their model
overestimates profit with reserve sales by 2 to 5 percentage points,
mainly because it lets units sell reserves below their minimum output.
Linear water values used in a mixed-integer simulation lose only about
one percentage point, which suggests that much of the linearisation
error arises in the operation within the week rather than in the water
values. \citet{abgottspon2014} use SDP between weeks together with hourly
subproblems within each week, and keep the small daily reservoir out of
the state by resetting it every week. \citet{perezdiaz2020medium} propose
SDP as a benchmark for quantifying the approximation errors of linear
policies.

The flexibility of SDP comes at a computational cost. The number of
storage-state combinations grows rapidly as reservoirs are added, and
the reviewed applications accordingly concern one plant or a small
cascade. Its suitability depends on the size of the state space as well
as on the benefit of retaining non-convex operating detail.

\subsubsection{Stochastic Dual Dynamic Programming}

Stochastic dual dynamic programming (SDDP), introduced by
\citet{pereira1991}, approximates future value functions using cuts,
avoiding a full grid over reservoir volumes, and is regarded as the
state-of-the-art method for medium- and long-term scheduling in
hydropower-dominated systems \citep{helseth2016}. The convergence
analysis of \citet{shapiro2011} concerns multistage linear stochastic
problems, and standard cut arguments do not extend directly to
non-convex formulations. Applications therefore approximate the physical
detail that SDP can retain. \citet{helseth2013} include pump start-up
costs through a continuous operating state instead of integer
variables. This reduces switching from generation to pumping sharply,
but the pump still runs below its intended minimum output, so the
approximation captures the effect of start-up costs without providing an
exact commitment schedule. In the same study, transmission constraints
and a finer, chronological time resolution change cable use and the
assessed profitability of the cable, and raise the run time several
times over. \citet{helseth2016} combine SDP for the price
process with SDDP for the reservoirs, which allows inflows, energy prices
and reserve prices to be uncertain, at a run time of more than two days,
and warn that the linear plant description may cause significant errors
when reserves are sold.

\paragraph{Extending the method}
Several studies trade computation time for physical accuracy.
\citet{hjelmeland2019energies} include head-dependent generation and use
strengthened cuts, which give a better operating policy at about five
times the run time, although the gap between the bounds remains large
when prices are uncertain. \citet{helseth2024approximating} connect an
aggregated representation of the Norwegian hydropower system with
feasibility cuts computed beforehand from the detailed system. This
lowers total cost compared with plain aggregation at almost three times
the run time, and is considerably faster than the higher-dimensional
cuts of the earlier method of \citet{helseth2023hydropowera}. Such
extensions narrow the gap between the two dynamic-programming approaches
without removing the underlying trade-off.

\paragraph{Applications to individual owners}
Applications of SDDP to an individual owner's valuation remain
concentrated on conventional reservoir hydropower.
\citet{bore2026assessing} use SDDP with weekly decision stages and a
linear generation--discharge relationship. Each iteration samples one
path through the weeks, the algorithm stops after a fixed number of
iterations, and no bounds are reported; they attribute a water value that
does not fall as the reservoir fills to the limited range of reservoir
levels visited. For pumped storage, \citet{abgottspon2016multihorizon}
use SDDP in a multi-horizon model but report a deterministic case, and
\citet{schuller2024practical} use water values from an SDDP-based model
to value a pumping investment in the energy market. Owner-oriented
pumped-storage studies that include reserve markets have relied on SDP
or deterministic optimisation.

The reviewed approaches consequently offer different combinations of
computational tractability, operational detail and adaptation to new
information. Two-stage formulations simplify the future information
structure; rolling-horizon planning repeatedly updates decisions;
discretised SDP can retain non-convex stage problems for small systems;
and SDDP exploits a convex structure to approximate future value for
larger ones. Their relative suitability follows from the decision
problem and the assumptions retained, rather than from a universal
ranking of the methods.

% ---------------------------------------------------------------------
\subsection{Synthesis and Research Gap}
\label{sec:lit_gap}

Table~\ref{tab:lit_comparison} summarises selected studies along the
dimensions discussed above. The literature establishes several relevant
strands: economic valuation from different perspectives, representation
of uncertainty, participation in multiple markets, comparison of plant
configurations, and approaches for solving the resulting scheduling
problems. These strands address complementary aspects of the value of
pumped-storage flexibility.

The closest studies cover different parts of the capacity-valuation
question, and each provides a different piece of evidence. A system
study shows that a large pumped-storage expansion has a lower marginal
value once the watercourse is represented in detail
\citep{helseth2017}. Owner studies show that adding a pump yields modest
gains in the energy market and that its value is limited by pumping
costs and short-term restrictions
\citep{ak2019quantifying,alic2023optimal,schuller2024practical}. Studies
of reserve markets show that the most valuable configuration is the one
with the widest operating range \citep{helseth2017iet,chazarra2017},
and an island study shows that reservoir size involves a trade-off
between cost and reliability \citep{zisos2023mixing}. These studies
differ in objectives, uncertainty, horizons and markets, as they were
designed to answer different questions. Their findings therefore cannot
be combined into a common assessment of pump, turbine and reservoir
capacities for a pumped-storage owner.

\citet{bore2026assessing} already connect stochastic multi-market
operation to the value of generating-capacity upgrades. The remaining
research need concerns the joint contribution of pumping, generating and
storage capacity in a pumped-storage setting. It also concerns the
influence of price dynamics: the reviewed studies relate price behaviour
to storage value only with known prices, and none separates volatility
from persistence, although the two may favour different capacities.

Against this background, the present study asks how pumping power,
generating power and reservoir capacity, individually and in
combination, affect the owner's long-term operating value across
electricity markets. A related question is how their relative value
changes with price volatility and persistence. The intended contribution
is a comparative understanding of these capacity interactions under
uncertainty, building on the scheduling and valuation approaches
established in the reviewed literature. The last row of
Table~\ref{tab:lit_comparison} places the study alongside the reviewed
literature.

% ---------------------------------------------------------------------
% Table 3. Check the row "This study" against ch. 6 when it changes
% (reserve products depend on Q-002).
\begin{landscape}
\begingroup
\footnotesize
\setlength{\tabcolsep}{4pt}
\renewcommand{\arraystretch}{1.15}
\begin{longtable}{@{}>{\raggedright\arraybackslash}p{0.12\linewidth}
                   >{\raggedright\arraybackslash}p{0.15\linewidth}
                   >{\raggedright\arraybackslash}p{0.18\linewidth}
                   >{\raggedright\arraybackslash}p{0.17\linewidth}
                   >{\raggedright\arraybackslash}p{0.17\linewidth}
                   >{\raggedright\arraybackslash}p{0.15\linewidth}@{}}
\caption{Selected studies compared by valuation perspective, modelling
approach, market participation and capacity question, and the
positioning of this study.}
\label{tab:lit_comparison}\\
\toprule
\textbf{Study} & \textbf{Perspective and asset} & \textbf{Time and uncertainty} & \textbf{Model representation} & \textbf{Market focus} & \textbf{Capacity question}\\
\midrule
\endfirsthead
\multicolumn{6}{l}{\tablename~\thetable\ (continued)}\\
\toprule
\textbf{Study} & \textbf{Perspective and asset} & \textbf{Time and uncertainty} & \textbf{Model representation} & \textbf{Market focus} & \textbf{Capacity question}\\
\midrule
\endhead
\midrule
\multicolumn{6}{r}{\footnotesize Continued on the next page}\\
\endfoot
\bottomrule
\addlinespace
\multicolumn{6}{@{}p{\linewidth}@{}}{\footnotesize SDP: stochastic
dynamic programming; SDDP: stochastic dual dynamic programming; NLP:
nonlinear programming; MILP: mixed-integer linear programming; PV:
photovoltaic; DA/ID: day-ahead/intraday; FCR: frequency containment
reserve (N: normal operation, D: disturbances); aFRR: automatic
frequency restoration reserve. System-perspective studies and general
methodological references are discussed in the text but are not
included in this focused comparison.}\\
\endlastfoot
\citet{abgottspon2013stochastic} & Price-making owner; Swiss pumped storage & One year; inflow & SDP with hourly recourse & DA and forward contracts & Operation of the studied plant\\
\citet{abgottspon2014} & Owner; Swiss pumped storage & One year; inflow and prices; known or uncertain week compared & SDP with stochastic or simplified within-week models & DA and weekly secondary reserve capacity & Operation of the studied plant\\
\citet{helseth2016} & Price-taking owner; reservoir hydro & Medium term; inflow, energy and reserve prices & Combined SDP/SDDP; linear & Sequential energy and reserve-capacity markets & Water values and schedules with reserve access\\
\citet{helseth2017iet} & Price-taking owner; reservoir hydro & One year; weekly states of inflow and prices & SDP with weekly MILP; head-dependent generation & DA energy and primary reserve capacity & Two reinvestment alternatives\\
\citet{chazarra2017} & Price-taking owner; closed-loop pumped storage & Daily scheduling; known prices and activation & MILP; hydraulic short-circuit operation & DA, secondary reserve and regulation energy & Three unit configurations\\
\citet{ak2019quantifying} & Owner; cascade with pumping & Annual evaluation; historical inflows and price years & NLP with monthly decisions and operating rules & DA energy & Existing cascade versus pumped-storage operation\\
\citet{alic2023optimal} & Owner; pumped storage in a cascade & One year; weekly states of inflow and energy prices & SDP; state-dependent environmental constraints & DA energy & Adding a \SI{38}{\mega\watt} pump; effect of constraints\\
\citet{zisos2023mixing} & Hybrid island-energy design & Long-run simulation; wind, demand and conversion & Rule-based simulation with heuristic search & Supply reliability and economic performance & Reservoir and PV sizing; fixed pump and turbine power\\
\citet{schuller2024practical} & Price-taking owner; Norwegian watercourse & Inflow and prices & SDDP-based water values with hourly simulation & Energy & A new plant with pumping\\
\citet{naval2025optimization} & Hybrid-system operator; pumped storage & Hourly operation over a year; wind; prices known & MILP with chance constraints & Grid exchange and renewable demand coverage & Operation of a specified hybrid configuration\\
\citet{bore2026assessing} & Price-taking owner; reservoir hydro & One year; weekly stages; inflow and prices; known within-week inputs & SDDP; linear generation model & DA, ID, FCR-N, FCR-D and aFRR capacity; activation excluded & Generating-power upgrades; maximum feasible investment cost\\
\midrule
This study & Price-taking owner; pumped storage with natural inflow & Weekly stages; inflow and DA prices; known within-week inputs; varied price volatility and persistence & SDDP; linear plant model & DA energy and reserve capacity; activation excluded & Combinations of pumping power, generating power and reservoir capacity\\
\end{longtable}
\endgroup
\end{landscape}
